\chapter{Large Language Model Tasks}
\label{ch:llm-tasks}

This chapter is the book's first sustained treatment of the engine's
AI--HPC integration layer, and the conceptual core of the workshop abstract
that motivates this textbook. \texttt{TaskType.LLM} invokes an
OpenAI-compatible Chat Completions API \emph{as a workflow task}: a Large
Language Model inference becomes a DAG node with the same dependency,
staging, and provenance guarantees as a numerical simulation.

\section{Design Goals}

\texttt{TaskType.LLM} is intended for JSON request construction, local
workflow data inputs, and response capture; notably, \textbf{it does not
require the OpenAI Python package}. The task speaks the OpenAI-compatible
Chat Completions HTTP surface directly, which means any provider exposing
that surface --- not only OpenAI itself --- can be configured as a target.

\section{Provider Configuration}

Each provider has an INI section named \texttt{llm.<provider-name>}; any
number of these sections may be configured simultaneously
(\Cref{sec:config-llm}). \texttt{endpoint} is the provider base URL (or a
full \texttt{/chat/completions} URL); \texttt{model} is the default model;
credentials come from either \texttt{api\_key\_env} or \texttt{api\_key} ---
the environment-variable form is preferred:

\begin{lstlisting}[style=ini]
[llm.research]
endpoint=https://api.example.org
api_key_env=DAGON_RESEARCH_LLM_KEY
model=research-chat
\end{lstlisting}

\begin{cautionbox}
Do not commit real API keys. The endpoint receives a \texttt{POST
/v1/chat/completions} request unless its configured URL already ends in
\texttt{/chat/completions}.
\end{cautionbox}

\section{Creating an LLM Task}

The prompt is a JSON-compatible Chat Completions request. \texttt{params}
supplies named values for \texttt{\{placeholder\}} strings anywhere in that
JSON. The \texttt{model} value, if present in the request, overrides the
provider default.

\begin{lstlisting}[caption={Minimal LLM task construction.}]
ask = DagonTask(
    TaskType.LLM,
    "summarize",
    {"messages": [{"role": "user", "content": "Summarize {subject}."}], "temperature": 0.2},
    provider="research",
    params={"subject": "the experiment"},
)
\end{lstlisting}

The prompt may instead be supplied as a JSON \emph{string}, which is useful
when a script assembles the request dynamically with \texttt{json.dumps}.
When a workflow is serialized, LLM tasks use \texttt{"type": "llm"} and
retain the provider, parameters, input-file mapping, and output filename;
loading that JSON elsewhere requires the named provider to be configured
locally on the loading machine.

\section{Workflow Input Files: Connecting Simulation Output to a Prompt}
\label{sec:llm-input-files}

\texttt{input\_files} attaches UTF-8 text produced by another task. Keys
become prompt parameters; values use the ordinary \texttt{workflow://} form
developed in \Cref{ch:schema}:

\begin{lstlisting}
input_files={"report": "workflow:///prepare/output/report.txt"}
\end{lstlisting}

This single line creates the \texttt{prepare} $\to$ \texttt{review}
dependency, exactly as in any file-based task. Before the request is sent,
the file is copied under
\texttt{.dagon/inputs/<workflow>/<task>/...} inside the LLM task's
directory, and its UTF-8 content replaces \texttt{\{report\}} in the prompt.
An inline \texttt{workflow:///prepare/output/report.txt} reference written
directly inside a prompt string is also parsed, staged, and replaced; using
\texttt{input\_files} is recommended chiefly because it keeps longer prompts
readable.

\begin{cautionbox}
Input paths must be relative to the producer task directory and cannot
contain \texttt{..}. Binary files and remote-only producer files are
\textbf{not} supported by this task type: the LLM task's file integration is
deliberately restricted to local, UTF-8 text.
\end{cautionbox}

The full response JSON is written to \texttt{response.json} in the LLM
task's working directory (or to the relative \texttt{output\_file} supplied
to the task), so a downstream task can consume it via \texttt{workflow://}
exactly as it would consume any other task's file output.

\section{Worked Example: A Locally Verifiable LLM Workflow}

The repository ships a complete, locally verifiable, credential-free example
that this book reproduces in full, because it is the cleanest illustration
of how an AI inference step is woven into an ordinary DAG.

\begin{lstlisting}[caption={\texttt{llm\_workflow.py}: a self-contained, mock-backed LLM workflow.}]
import json
import tempfile
import threading
from http.server import BaseHTTPRequestHandler, ThreadingHTTPServer
from pathlib import Path

from dagon import Workflow
from dagon.task import DagonTask, TaskType


class MockChatHandler(BaseHTTPRequestHandler):
    def do_POST(self):
        if self.path != "/v1/chat/completions":
            self.send_error(404)
            return
        request = json.loads(self.rfile.read(int(self.headers["Content-Length"])).decode("utf-8"))
        content = request["messages"][-1]["content"]
        payload = json.dumps({"choices": [{"message": {"role": "assistant", "content": "Mock summary: " + content}}]}).encode("utf-8")
        self.send_response(200)
        self.send_header("Content-Type", "application/json")
        self.send_header("Content-Length", str(len(payload)))
        self.end_headers()
        self.wfile.write(payload)

    def log_message(self, _format, *_args):
        pass


def main():
    server = ThreadingHTTPServer(("127.0.0.1", 0), MockChatHandler)
    threading.Thread(target=server.serve_forever, daemon=True).start()
    try:
        with tempfile.TemporaryDirectory(prefix="dagon-llm-lesson-") as scratch:
            config = {
                "batch": {"scratch_dir_base": scratch, "run_base": "", "threads": "1"},
                "dagon_service": {"use": "False", "route": "http://localhost:57000"},
                "ftp_pub": {"ip": "localhost"}, "slurm": {"partition": ""},
                "llm.local_mock": {"endpoint": "http://127.0.0.1:%s" % server.server_port,
                                   "api_key": "example-key-not-a-secret", "model": "mock-chat"},
            }
            workflow = Workflow("Lesson12", config=config)
            produce = DagonTask(TaskType.BATCH, "prepare_report", "mkdir -p output; printf 'temperature increased by 1.2 C' > output/report.txt")
            summarize = DagonTask(TaskType.LLM, "summarize_report",
                {"messages": [{"role": "user", "content": "Summarize this report: {report}"}]},
                provider="local_mock", input_files={"report": "workflow:///prepare_report/output/report.txt"})
            workflow.add_task(produce)
            workflow.add_task(summarize)
            workflow.run()
            print(Path(summarize.working_dir, "response.json").read_text(encoding="utf-8"))
    finally:
        server.shutdown()
        server.server_close()


if __name__ == "__main__":
    main()
\end{lstlisting}

\begin{lstlisting}[style=bash]
python3 llm_workflow.py   # the listing above, saved to a file
\end{lstlisting}

The listing runs unchanged against DAGonStar~3.0.0. The course versions of
the same mechanism are \texttt{examples/27-llm/llm\_task.py}
(Lecture~27: templates, staged inputs, providers) and
\texttt{examples/28-llm-audit/llm\_audit.py} (Lecture~28: what is recorded,
and rehearsal with portable emulation).

\begin{figure}[h]
\centering
\begin{tikzpicture}[
  node distance=20mm,
  task/.style={rectangle, rounded corners, draw=dagonblue, very thick,
    fill=blue!4, minimum width=34mm, minimum height=12mm, font=\sffamily\small, align=center},
  ext/.style={rectangle, rounded corners, draw=dagonaccent, very thick,
    fill=teal!6, minimum width=34mm, minimum height=12mm, font=\sffamily\small, align=center},
  arr/.style={-{Stealth[length=2.6mm]}, thick, dagonblue}
]
  \node[task] (prep) {prepare\_report\\ \footnotesize writes \texttt{output/report.txt}};
  \node[task, right=28mm of prep] (llm) {summarize\_report\\ \footnotesize \texttt{TaskType.LLM}};
  \node[ext, right=22mm of llm] (api) {OpenAI-compatible\\ \footnotesize Chat Completions API};
  \draw[arr] (prep) -- node[above,font=\sffamily\scriptsize]{\texttt{input\_files=\{report: workflow://...\}}} (llm);
  \draw[arr] (llm) -- node[above,font=\sffamily\scriptsize]{POST /v1/chat/completions} (api);
  \draw[arr] (api.south) .. controls +(0,-1.2) and +(2.2,-1.2) .. node[below,font=\sffamily\scriptsize]{JSON response} (llm.south);
  \node[task, below=14mm of llm, text width=34mm] (resp) {\texttt{response.json}\\ \footnotesize available to downstream tasks via \texttt{workflow://}};
  \draw[arr] (llm) -- (resp);
\end{tikzpicture}
\caption{An LLM task as a first-class DAG node: it consumes a producer's text
output through the ordinary \texttt{workflow://} mechanism and exposes its
JSON response as an ordinary, dependency-trackable task output.}
\label{fig:llm-task-flow}
\end{figure}

\section{Walkthrough}

\subsection{Step 1: Execute and Observe}

The mock server implements only the minimal OpenAI-compatible
\texttt{/v1/chat/completions} surface needed for the lesson, bound to
\texttt{127.0.0.1} on an unused port, and is shut down automatically when the
workflow completes. Expected output is a JSON completion of the shape:

\begin{lstlisting}
{"choices": [{"message": {"content": "Mock summary: Summarize this report: temperature increased by 1.2 C"}}]}
\end{lstlisting}

\subsection{Step 2: The Producer Task}

\begin{lstlisting}
produce = DagonTask(
    TaskType.BATCH,
    "prepare_report",
    "mkdir -p output; printf 'temperature increased by 1.2 C' > output/report.txt",
)
\end{lstlisting}

Its eventual file location is conceptually
\texttt{<prepare\_report task directory>/output/report.txt} --- an absolute
path generated only at runtime, and therefore never written directly into
the LLM task's prompt, which would make the workflow non-portable.

\subsection{Step 3: The JSON Request and File Attachment}

\begin{lstlisting}
summarize = DagonTask(
    TaskType.LLM,
    "summarize_report",
    {"messages": [{"role": "user", "content": "Summarize this report: {report}"}]},
    provider="local_mock",
    input_files={"report": "workflow:///prepare_report/output/report.txt"},
)
\end{lstlisting}

\texttt{\{report\}} is a prompt parameter; \texttt{input\_files} supplies its
value via the standard \texttt{workflow:///<producer-task>/<path>} form.
Building dependencies recognizes the reference and creates
\texttt{prepare\_report} $\to$ \texttt{summarize\_report}, exactly as
\Cref{ch:schema} predicts for any other task type.

\subsection{Step 4: Provider Configuration in Code}

Because the mock server selects a free port at runtime, the example builds
its provider configuration in Python rather than in a static INI file:

\begin{lstlisting}
"llm.local_mock": {
    "endpoint": "http://127.0.0.1:%s" % server.server_port,
    "api_key": "example-key-not-a-secret",
    "model": "mock-chat",
}
\end{lstlisting}

Provider section names always use the \texttt{llm.<provider-name>} pattern;
the task's \texttt{provider} argument selects the matching suffix.

\subsection{Step 5: Run and Inspect}

\begin{lstlisting}
workflow.add_task(produce)
workflow.add_task(summarize)
workflow.run()
\end{lstlisting}

\texttt{run()} builds missing dependencies automatically. On completion,
\texttt{summarize.working\_dir} identifies the LLM task's directory, and the
full response is available at \texttt{Path(summarize.working\_dir, "response.json")}.
The \emph{complete} JSON response is deliberately retained, rather than only
the assistant's text, so a downstream task can extract provider-specific
metadata or choices with its own deterministic program --- a design choice
directly in service of provenance.

\subsection{Step 6: Switching to a Real Provider}

For a real OpenAI-compatible service, the workflow request is unchanged; only
configuration changes, in a local, ignored \texttt{dagon.ini}:

\begin{lstlisting}[style=ini]
[llm.research]
endpoint=https://api.example.org
api_key_env=DAGON_RESEARCH_LLM_KEY
model=research-chat
\end{lstlisting}

\begin{lstlisting}[style=bash]
export DAGON_RESEARCH_LLM_KEY=...
\end{lstlisting}

with \texttt{provider="research"} substituted into the task constructor.
\begin{cautionbox}
Never commit the key, and never replace the safe sample configuration with a
real one in version control.
\end{cautionbox}

\section{Verification}

\begin{lstlisting}[style=bash]
python3 -m unittest tests.test_llm -v
\end{lstlisting}

\begin{practicebox}
Record the provider name, model, prompt template, and source input files in
your workflow source, and preserve \texttt{response.json} as a workflow
output. This makes AI-assisted interpretation inspectable --- a downstream
reader can audit exactly what was asked and what was returned --- while
keeping secret credentials entirely out of the workflow and the repository.
\end{practicebox}

\section{Why a Provider-Agnostic Surface Matters}

\Cref{ch:llm-tasks} already noted that \texttt{TaskType.LLM} requires no
provider-specific SDK, speaking the OpenAI-compatible Chat Completions
surface directly over HTTP. This design decision is worth examining for its
broader methodological implications, because it bears directly on the
scientific reproducibility properties this book has emphasized throughout.

A workflow engine that depended on a specific vendor's Python SDK would
couple every LLM-using workflow to that SDK's release cadence, versioning
behavior, and continued maintenance --- a dependency risk qualitatively
different from, and arguably larger than, the dependency risk already
accepted for, say, the Docker SDK (\Cref{ch:backends}), since the
generative-AI provider landscape has historically changed faster than
container runtimes have. By speaking the request/response \emph{protocol}
directly, rather than a vendor SDK, \texttt{TaskType.LLM} decouples a
workflow's AI-assisted step from any single vendor's continued existence:
the \texttt{endpoint} configuration value (\Cref{sec:config-llm}) is, in
principle, swappable to any provider implementing the same
\texttt{/chat/completions} surface, with no change to the workflow's Python
source.

\begin{practicebox}
This protocol-level portability is the same architectural strategy already
examined, in a different guise, for remote execution
(\Cref{sec:factory-pattern-depth}): just as a \texttt{BATCH} task's command
string is portable across any POSIX-like remote host reachable by SSH, an
\texttt{LLM} task's JSON request is portable across any provider speaking
the OpenAI-compatible surface, with the concrete provider selected entirely
by configuration (\Cref{ch:configuration}), not by code.
\end{practicebox}

\section{The Epistemics of AI-Assisted Scientific Interpretation}

Embedding an LLM call inside a scientific DAG, as Lecture~27 demonstrates,
raises a methodological question that is worth addressing directly, rather
than leaving implicit: what kind of scientific claim, if any, does the
output of an \texttt{LLM} task support, and how should a researcher's
confidence in that output be calibrated?

It is important to be precise about what the mock example of
\Cref{ch:llm-tasks} is, and is not, demonstrating. The mechanism --- a
prompt assembled from a producer task's UTF-8 output, sent to a configured
endpoint, with the full response preserved as a dependency-trackable file
--- is provider-agnostic and equally applicable whether the configured
\texttt{provider} ultimately resolves to a deterministic mock, a
well-characterized classical text-processing service, or a state-of-the-art
generative language model. The engine itself makes no claim about, and has
no mechanism for assessing, the \emph{factual reliability} of whatever text
a given provider returns; that assessment is and remains entirely the
responsibility of the researcher consuming the response.

\begin{cautionbox}
A workflow that uses an \texttt{LLM} task to \emph{summarize} or
\emph{describe} a model's numerical output (as in the worked example of
this chapter) is making a categorically different kind of use than a
workflow that uses an \texttt{LLM} task's output as a direct, unreviewed
input to a further numerical or decision-making step. The former is
analogous to using a research assistant's draft summary --- valuable as a
starting point, properly understood as requiring expert review before being
relied upon; the latter would require a substantially higher standard of
verification, calibration, and, in most scientific contexts, human review
in the loop. This book's worked examples are deliberately confined to the
former, lower-stakes category, and the full-response-retention design
discussed below is precisely what makes the necessary expert review
possible after the fact.
\end{cautionbox}

This is, in fact, the deeper reason \Cref{ch:llm-tasks} emphasized that the
\emph{complete} response JSON, not merely the assistant's extracted text, is
preserved in \texttt{response.json}: an unreviewed text summary, divorced
from the full request that produced it (including the exact prompt,
provider, and model), is far harder for a later reader to critically assess
than a complete, inspectable request/response pair. The provenance argument
already developed in \Cref{ch:schema} for file-based dependencies applies
with, if anything, greater force here: the epistemic status of an
AI-generated scientific summary depends critically on knowing exactly what
was asked and exactly what model produced the answer, information that a
bare text string alone does not carry.

\section{Cost, Latency, and Rate Limits as Operational Constraints}

A final consideration, absent from the purely mechanical treatment of
Lecture~27 but essential for any researcher moving from the mock example to a
real provider, concerns the operational constraints that distinguish LLM
tasks from the local, cost-free \texttt{BATCH} tasks used throughout most of
this book's tutorial examples.

Unlike a local shell command, a real-provider LLM completion typically
incurs a per-request monetary cost, a network round-trip latency
significantly larger than a local subprocess call, and is frequently
subject to rate limits imposed by the provider. These constraints interact
directly with the workflow-design questions already raised elsewhere in
this book. The task-granularity discussion of \Cref{sec:task-granularity},
for instance, takes on a sharper edge for LLM tasks specifically: an
overly fine-grained decomposition that issues many small, separate LLM
requests where one consolidated request would suffice multiplies both cost
and the risk of rate-limit-induced failures, with no corresponding gain in
traceability beyond what a single, well-structured request and response
already provide.

\begin{practicebox}
The \texttt{timeout} parameter exposed on the \texttt{LLMTask} constructor
(Appendix~A, default 120 seconds) is the engine's primary mechanism for
bounding the operational risk of an unresponsive or rate-limited provider;
a workflow author moving from the mock provider used in this book's local
examples to a production LLM provider should explicitly reconsider this
default against the provider's documented typical latency and the
workflow's own tolerance for a stalled task, rather than relying on the
default value chosen for local, low-latency testing.
\end{practicebox}

\section{Chapter Summary}

\texttt{TaskType.LLM} demonstrates the engine's central architectural claim
at its most consequential: an AI inference call, expressed as a JSON Chat
Completions request, participates in the same dependency graph, the same
\texttt{workflow://} staging mechanism, and the same provenance discipline as
a shell command or an HPC job submission. This chapter additionally argued
that the task type's provider-agnostic, protocol-level design is itself a
reproducibility decision rather than a mere implementation convenience,
addressed directly the epistemic status of AI-generated scientific text
embedded in a DAG and the corresponding case for full-response retention,
and discussed the cost, latency, and rate-limit constraints that distinguish
LLM tasks operationally from the local tasks examined in Part~I.
\Cref{ch:web-tasks} extends the same architectural principle to arbitrary
HTTP/HTTPS services, and \Cref{ch:native-tasks} extends it to native Python
functions.

\begin{exercise}
A researcher proposes replacing a workflow's \texttt{NATIVE}-task
statistical quality-control step (\Cref{ch:native-tasks}) with an
\texttt{LLM} task that asks a language model to ``check whether this data
looks reasonable.'' Using the epistemic-status discussion of this chapter,
write a short methodological critique of this proposal, and suggest a
revised design that uses an \texttt{LLM} task appropriately alongside,
rather than instead of, a deterministic quality-control mechanism.
\end{exercise}
